% concatenated from NilpotentPrimitive.tex
\documentclass[12pt]{amsart}
\usepackage{amssymb,euscript}
\setlength\textwidth{36pc}
\setlength\textheight{53pc} \setlength\oddsidemargin{16pt}
\setlength\evensidemargin{16pt}
%\renewcommand{\baselinestretch}{1.05}
\raggedbottom
\relpenalty=10000
\binoppenalty=10000
\tolerance=500
\mathsurround=1pt
\newtheorem{theorem}{Theorem}[section]
\newtheorem{lemma}[theorem]{Lemma}
\newtheorem{proposition}[theorem]{Proposition}
\newtheorem{corollary}[theorem]{Corollary}
\theoremstyle{definition}
\newtheorem{definition}{Definition}
\newtheorem{example}{Example}
\theoremstyle{remark}
\newtheorem*{remark}{Remark}
\newtheorem*{notation}{Notation}
\renewcommand{\hat}{\widehat}
\renewcommand{\tilde}{\widetilde}

%This is the Comm. Algebra paper, difficult to access through
%journal website. Checked against published version; matches
%published version.
%On DF w^3 google homepage; on arXiv.

%----------

%Incorporating revisions from detinko e-mail 8 August 2003

%Revisions 14/10/03 (stupid error)

\begin{document}

%\bigskip

%\bigskip

\title{Nilpotent primitive linear groups over finite fields}

%\bigskip

%\bigskip

\author{A. S. Detinko}
\address{ Department of Computer Science,
School of Computing and Engineering,
University of Huddersfield, Queensgate,
Huddersfield
HD1 3DH, UK}
\email{detinko.alla@gmail.com}
\author{D. L. Flannery}
\address{Department of Mathematics, University
 of Galway, Galway H91 TK33, Ireland}
\email{dane.flannery@universityofgalway.ie}
%\subjclass{Primary 20; Secondary 20}

\maketitle

%\bigskip

%\vspace*{20mm}

\section{Introduction}

 In this paper we investigate the structure
 of groups as in the title. Our work builds on work of
 several other authors, namely Konyukh \cite{Konyuh},
 Leedham-Green and Plesken \cite{LGP}, 
 and  Zalesskii \cite{Zalesskii}, who  have described 
 the abstract isomorphism types of the groups. We obtain 
 more detailed descriptions, in particular
 explaining how group structure depends on the existence 
 of an abelian primitive 
%maximal abelian 
 subgroup. Additionally we show that isomorphism
 type of a group completely determines its 
 conjugacy class in the relevant general linear group.

 A  brief outline of the paper now follows.
 In Section~\ref{ablin} we review standard material on
 abelian (cyclic) irreducible linear groups.
 In Section~\ref{fund} fundamental structural results are 
 given. In Section~\ref{deg2} nilpotent primitive linear 
 groups of degree 2 are classified up to conjugacy, 
 and then groups of degree greater than 2 are treated 
 thoroughly in Section~\ref{greater2}. 
 The final Section~\ref{summary} summarises our results. 

 Throughout, $\mathbb{F}$ is a finite field of size $q$ and 
 characteristic $p$, and $G \leq \mathrm{GL}(n,\mathbb{F})$,
 $n>1$. The natural (right) $\mathbb{F}G$-module 
 of dimension $n$ is denoted $V$. Whenever we
 refer to  ``primitive'' or ``imprimitive'' 
 linear groups, we are implicitly assuming 
 them to be irreducible.
%
% Although we focus on linear groups over finite 
% fields, some results pertain as well to groups over other 
% fields of positive or even zero characteristic.

\section{Abelian linear groups}
\label{ablin}

 %Most of the theory in this section is quite well-known.
 Suppose $A$ is an abelian irreducible subgroup of 
 $\mathrm{\mathrm{GL}}(n, \mathbb{F})$, and   
 denote by $\Delta$ the enveloping algebra 
 $\langle A  \rangle _\mathbb{F}$ of $A$.
%That is, $\Delta$ is the smallest 
% $\mathbb{F}$-subalgebra of the full matrix algebra 
% $\mathrm{Mat}(n,\mathbb{F})$ containing $A$. 
\begin{proposition}
\label{singdef}
 $\Delta$ is a field, $|\Delta : \mathbb{F}1_n| =n$,
 so that $A$ is cyclic of order dividing $q^n-1$. 
 Hence  $\Delta^{\! \times}$ is a maximal abelian  
 subgroup of $\mathrm{GL}(n, \mathbb{F})$.
\end{proposition}
\begin{proof}
 The first two claims are covered by 
 \cite[Theorem 4, p.99]{Suprunenko}. The rest is then clear.
 %Suppose $B$ is an abelian subgroup of 
 %$\mathrm{GL}(n, \mathbb{F})$ containing $\Delta^{\! \times}$. 
 %By the above, $\langle B \rangle_{\mathbb{F}}$ has order 
 %$q^n$ $=$ $|\Delta|$, and therefore $B$ $=$ $\Delta^{\! \times}$. 
\end{proof}
\begin{corollary}
 A subgroup $C$ of $\Delta^{\! \times}$ is irreducible if 
 and only if
 $|\langle C \rangle_{\mathbb{F}} : \mathbb{F}1_n| =n$.
\end{corollary} 
%\begin{lemma}
%\label{singcyc}
% There exist elements of ${\rm GL}(n,\mathbb{F})$ of order
% $q^n-1$, and if $g$ is any such element then
% $\langle g  \rangle$ is irreducible.
%\end{lemma}
%\begin{proof}
% Let $\alpha$ be a generator of $\mathrm{GF}(q^n)^{\times}$.
% Then 
% $\mathrm{diag}(\alpha, \alpha^q, \ldots \, \! , \alpha^{q^{n-1}})$
% has order $q^n-1$ and trace in $\mathbb{F}$, so is 
% conjugate to an element of $\mathrm{GL}(n,\mathbb{F})$
% by \cite[(9.5)(c), p.147]{Isaacs}
% and \cite[(9.23), p.155]{Isaacs}.
% 
% Let $g \in \mathrm{GL}(n,\mathbb{F})$,
% $|g| = q^n-1$.
% By Maschke's Theorem,  $G = \langle g  \rangle$ is 
% completely reducible. 
% Suppose  $G$ is reducible, with $m$ irreducible parts, 
% each of degree $n_i$, $1< i \leq m$. Then $|G|$ divides 
% $r=(q^{n_1}-1) \cdots (q^{n_m}-1)$ by Proposition~\ref{singdef}; 
% but $r\leq q^n-2$ by induction.  
%\end{proof}

 There always exist elements of ${\rm GL}(n,\mathbb{F})$ of 
 order $q^n-1$, and if $g$ is any such element then
 $\langle g  \rangle$ is irreducible.
 A cyclic subgroup of $\mathrm{GL}(n,\mathbb{F})$ of order
 $q^n-1$ is called a Singer cycle. The next few results
% are folklore. 
 about Singer cycles are well-known.
\begin{proposition}
\label{charirredab}
 A subgroup $C$ of $\Delta^{\! \times}$ is irreducible if 
 and only if $|C|$ does not divide $q^m-1$ for any 
 proper divisor $m$ of $n$.  
\end{proposition}
\begin{proposition}
\label{normcentr}
 Let $C$ be an irreducible subgroup of $\Delta^{\! \times}$.
 The $\mathrm{GL}(n,\mathbb{F})$-centraliser of 
 $C$ is $\Delta^{\! \times}$, and the 
 $\mathrm{GL}(n,\mathbb{F})$-normaliser of 
 $C$ is the semidirect product of   
 $\Delta^{\! \times}$ with a cyclic group of order
 $n$.
\end{proposition}
%suppose c=|C| does not divide q^m-1 for any divisor m of
%n, but divides q^d-1 for d<n not dividing n. Then c divides 
%q^d(q^m-d-1), so by induction divides q^d'-1for d'<d.
%By induction c divides q-1  
\begin{proposition}
\label{conjab}
 %All of the maximal irreducible abelian subgroups of 
 %$\mathrm{GL}(n,\mathbb{F})$ are 
 %conjugate to each other.
 %Thus 
 Abelian irreducible subgroups 
 of $\mathrm{GL}(n,\mathbb{F})$
 of the same order are conjugate.
\end{proposition}

 Now suppose $A$ is imprimitive, with imprimitivity
 system $\mathcal{S}$ $=$ $\{ V_1, \ldots \! \, ,V_k\}$, 
 $k > 1$. That is, the 
 components $V_i$ are subspaces of $V$,  
 $V = V_1 \oplus \cdots \oplus V_k$, 
 and $A$ permutes the $V_i$ by its natural action.
 In fact $A$ acts transitively because it is irreducible. 
 Hence $k$ divides $|A|$ and $n$, and 
%since $|A|$ divides $q^n-1$, $k$ and $q$ are coprime.
 $k$ is not divisible by $p$.
 
 Let $\phi$ be the homomorphism  $A \rightarrow \mathrm{Sym}(k)$
 arising from the
% permutation 
 action of $A$ on $\mathcal{S}$. An abelian transitive permutation 
 group is regular, so $\phi(A)$ is cyclic of order $k$.     
\begin{lemma}
\label{sizeprim}
 For each integer $r>1$ dividing $k$, $A$  
 has a system of imprimitivity consisting of $r$
 components.
\end{lemma}
\begin{proof}
 Set $k/r=l$ and $\phi(A) = \langle h\rangle$.  
 Then $\mathcal{S}$ is a disjoint union of $r$ 
 $\langle h^r \rangle$-orbits 
 $\{ V_{i_1} , \ldots \, \! , V_{i_l}\}$ of length $l$, 
 %Then  By \cite[Theorem 4, p.13]{Suprunenko},
 %it is readily checked that 
 %$\{ U_1 , \ldots \, \! ,  U_r\}$ is an 
 and $\{ V_{i_1} \oplus \cdots \oplus V_{i_l} 
 \mid 1 \leq i \leq r \}$ 
 is an $A$-system of imprimitivity.
\end{proof}

 By Lemma~\ref{sizeprim} we can consider that 
 $A$ has an imprimitivity system of prime size.
\begin{proposition}
\label{likesim}
 Let $k$ be prime, and denote $\ker \phi$ by $A'$.
\begin{itemize}
\item[{\rm (i)}]  Each $V_i$   
 is a faithful irreducible $A'$-module.
\item[{\rm (ii)}]  
 $|A|$ divides $k(q^{n/k}-1)$.
\end{itemize}
\end{proposition}
\begin{proof}
 The stabiliser $A_i$ $=$
 $\{ a\in A \mid V_i a =V_i \}$ acts irreducibly
 on $V_i$, $1\leq i \leq k$  
 (by \mbox{e.g.} \cite[Theorem 4.2B, p.68]{Dixon}).
 We have $A'= \cap_i A_i \leq A_i$, and because $|A : A'|$ is prime
 and $A$ is irreducible, $A'=A_i$ for all $i$. 
 If $a \in A$ fixes $V_i$ elementwise then by transitivity
 of $A$ on $\mathcal{S}$, $a$ fixes every element of
 $V$. That is, $V_i$ is a faithful $A_i$-module. Hence (ii) 
 follows by Proposition~\ref{singdef}.   
\end{proof}

 The converse of Proposition~\ref{likesim} (ii) is also true
 (see \cite[2.6]{Sim}).
\begin{proposition}
\label{likersim}
 A subgroup $C$ of  $\Delta^{\! \times}$  is 
 imprimitive if and only if $|C|$ divides 
 $k(q^{n/k}-1)$ for some prime $k$ dividing  $n$.
\end{proposition}

\section{Fundamental structural results}
\label{fund}

 We frequently use the following two lemmas, 
 whose straightforward proofs are omitted.
\begin{lemma}
\label{abnorsgp}
 A finite abelian normal subgroup of a primitive
 linear group 
 {\em (}over an arbitrary field\hspace*{.7mm}{\em )} 
 is cyclic.
\end{lemma}

\begin{lemma}
\label{ind2prim} 
 Any subgroup of index $2$ of a primitive linear group 
 is irreducible.
\end{lemma}

\begin{notation}
 If $H$ is nilpotent then we denote its Sylow $2$-subgroup 
 by $H_2$.
\end{notation}

 The following theorem 
%gives basic structural information. It 
 is vital.
\begin{theorem}
\label{konzal}
 Suppose $G$ is nilpotent primitive.  
\begin{itemize}
%\item[{\rm (i)}] $[G,G]$ is cyclic,
\item[{\rm (i)}] Every odd order subgroup of $G$ is cyclic.
\item[{\rm (ii)}] If $G_2$ is nonabelian then 
 $q\equiv 3$ mod $4$, and $G_2$ has maximal nilpotency 
 class: either $G_2 \cong Q_8$,
 or $G_2$ has order at least $16$ 
 and is generalised quaternion, 
 dihedral, or semidihedral. 
\end{itemize}
\end{theorem}
\begin{proof}
 See \cite{Zalesskii}
%.
 or \cite{Konyuh}. 
%All statements can 
 Everything  derives from the fact that $[G,G]$ is cyclic.
%; cf. \cite{Konyuh}.
% For (i) see \cite{Konyuh}.  The other claims
% can then be proved by straightforward means---or 
% see \cite{Zalesskii}. 
%(Parts (ii) and (iii) can be proved by straightforward 
%means, given (i). To see that (iii) implies $q \equiv 3$ 
%mod 4, note that if $q \equiv 1$ mod 4 then we can use 
%a scalar of order 4 to construct noncyclic abelian 
%normal subgroups of $G$,
% contradicting Lemma~\ref{abnorsgp}.)
\end{proof}
 
 Thus nilpotent primitive $G$ has the form $G_2 \times C$, 
 where the cyclic group $C$ is the direct product of the odd 
 order Sylow subgroups of $G$, and $G$ has a cyclic subgroup 
 of index 2. 

 A nilpotent primitive linear group over a finite field is 
 metacyclic (normalises a Singer cycle). Sim \cite{Sim} 
 classifies all metacyclic primitive linear groups of odd 
 prime power degree over finite fields. That case is not 
 relevant to our purposes, by the next proposition. 
\begin{proposition}
\label{oddfact}
 If $G$ is a  nonabelian nilpotent primitive subgroup
 of ${\rm GL}(n , \mathbb{F})$ then $n$ is even.
\end{proposition}
\begin{proof}
 By Theorem~\ref{konzal}, $q\equiv 3$ mod 4
 and $G$ has a cyclic normal subgroup $H$ 
 of index 2 and order divisible by 4. 
 By Lemma~\ref{ind2prim} and Proposition~\ref{singdef}, 
 $|H|$ divides $q^n-1$, whereas if $n$ is odd 
 then $q^n-1$ is not divisible by 4.
\end{proof}


\section{Degree 2}
\label{deg2}

\begin{proposition}
\label{absirredcase1}
 Suppose $G$ is nilpotent primitive. 
 Then $G$ is absolutely irreducible if and only 
 if $n=2$ and $G$ is nonabelian. 
\end{proposition}
\begin{proof} 
 Since every element of $G$ is semisimple  
 (\cite[Corollary 1, p.239]{Suprunenko}),
 $p$ does not divide $|G|$. 
 Thus $G$ is completely reducible over the algebraic 
 closure of $\mathbb{F}$, so 
 that if $n=2$ and $G$ is nonabelian then $G$ must be 
 absolutely  irreducible.
 Conversely, if $G$ is absolutely irreducible then it 
 is nonabelian, and isomorphic to an irreducible subgroup 
 of $\mathrm{GL}(n,\mathbb{C})$---see e.g. 
 \cite[Corollary 3.8, p.62]{Dixon}.
 By Theorem~\ref{konzal}, 
 $G$ has an abelian subgroup of index 2.  
 Therefore $n$ divides 2 by Ito's result 
 \cite[6.15, p.84]{Isaacs}.
\end{proof}

\begin{proposition}
\label{secchar}
 A nonabelian subgroup $G$ of
 $\mathrm{GL}(2,\mathbb{F})$ is nilpotent primitive 
 if and only if $q \equiv 3$ mod $4$ 
 and $G = G_2 \times C$, where $G_2$ is primitive
 and $C$ is a {\rm (}cyclic{\rm )} group of odd order 
 scalars.
\end{proposition}
\begin{proof}
 Suppose $G$ is nilpotent primitive, so $q \equiv 3$ mod 4.
 Then a Sylow 2-subgroup of the full monomial subgroup 
 of $\mathrm{GL}(2,\mathbb{F})$ is dihedral of order 8. 
 By Theorem~\ref{konzal} (ii), $G_2$ must be primitive, 
 because it is not monomial (the only sort of imprimitivity 
 in prime degree). Furthermore, $G_2$ is absolutely 
 irreducible by Proposition~\ref{absirredcase1}. Schur's 
 Lemma implies $C$ is scalar. 
\end{proof}
\begin{remark}
 We emphasise that the isomorphism types for $G_2$ 
 in Proposition~\ref{secchar} may be written down 
 from Theorem~\ref{konzal} (ii); cf. also \cite[II.2]{LGP}.
\end{remark}
\begin{remark}
 By \cite[III.4 (ii)]{LGP}, when $q\equiv 3$ mod 4,
 a nilpotent primitive subgroup of $\mathrm{GL}(n,\mathbb{F})$
 can have primitive Sylow 2-subgroup only if $n=2$.
 %Indeed, a 2-subgroup of $\mathrm{GL}(n,\mathbb{F})$ is 
 %reducible if $n>2$. ?? FALSE!! WHAT was the argument?
\end{remark}

 Below we say a little more about 2-subgroups of 
 $\mathrm{GL}(2,\mathbb{F})$. Prior to that we present
 some results on conjugacy between nilpotent linear groups.
\begin{theorem}
\label{isoeqconj}
 Suppose $n\neq p$ is prime, and 
 $G$, $H$ are nonabelian $n$-subgroups of
 $\mathrm{GL}(n,\mathbb{F})$. Then $G$, $H$ are conjugate 
 in $\mathrm{GL}(n,\mathbb{F})$ if and only if 
 $G \cong H$.
\end{theorem}
\begin{proof}
 If $G \cong H$ then by \cite[Proposition 4.2]{Conlon},
 $G$ and $H$ are conjugate as subgroups of 
 $\mathrm{GL}(n,\overline{\mathbb{F}})$, where
 $\overline{\mathbb{F}}$ is the algebraic closure of
 $\mathbb{F}$. 
% (Note that any completely reducible nilpotent subgroup of 
% $\mathrm{GL}(n,\overline{\mathbb{F}})$ is conjugate
% to a group of monomial matrices over $\overline{\mathbb{F}}$.)
 Thus $G$ and $H$ are conjugate in $\mathrm{GL}(n,\mathbb{F})$ 
 by the Deuring-Noether Theorem 
 \cite[1.22, p.26]{HuppertBlackburn}.
\end{proof} 
\begin{corollary}
\label{conl}
 Two nilpotent primitive subgroups of 
 $\mathrm{GL}(2,\mathbb{F})$ are
 conjugate in $\mathrm{GL}(2,\mathbb{F})$ if and only
 if they are isomorphic.
\end{corollary}
\begin{proof}
 The result is clear for abelian groups by
 Proposition~\ref{conjab}; otherwise  
 we appeal to Proposition~\ref{secchar} and 
 Theorem~\ref{isoeqconj}.
\end{proof}

\begin{lemma}
\label{2gps}
 Suppose $q\equiv 3$ mod $4$ and let $\omega$ be 
 a generator of ${\rm GF}(q^2)^{\times}$. 
 If $2^t$ is the largest power of $2$ dividing 
 $q+1$ then the subgroup of  
 ${\rm GL}(2,\mathbb{F})$  generated by 
$$
x= \left(
\begin{array}{cc}
 0 & 1 \\
 -\omega^{q+1} & \omega + \omega^q 
\end{array} \right)^{(q^2-1)/2^{t+1}}
 , \ \ \ \ \ \ \ 
y= \left(
\begin{array}{cc}
 1 & 0 \\
 \omega + \omega^q & -1 
\end{array}\right) 
$$
 
\noindent is a Sylow $2$-subgroup, and is semidihedral of
 order $2^{t+2}$.
 There are dihedral and generalised quaternion subgroups
 of ${\rm GL}(2,\mathbb{F})$ of every order $2^s$, 
 $3 \leq s \leq t+1$, unique up to conjugacy 
 at every order. Up to conjugacy, the unique 
 {\rm (}nonabelian{\rm )} semidihedral $2$-group in  
 ${\rm GL}(2,\mathbb{F})$ is the Sylow $2$-subgroup 
 indicated above.
\end{lemma}
\begin{proof}
% Note that $2^{a+1}$ is the largest power of
% 2 dividing $q^2-1$. 
 We readily verify that $|x| = 2^{t+1}$, $|y|=2$, and 
 $T = \langle x \rangle \rtimes \langle  y \rangle$
 is semidihedral. 
% Thus $T$ is a Sylow 2-subgroup of $\mathrm{GL}(2,\mathbb{F})$  
% by \cite[III.4]{LGP}. OR SIMPLY JUST COMPUTE ORDERS!! 
% Let $D$ be a semidihedral subgroup of 
% $\mathrm{GL}(2, \mathbb{F})$; then
% $|D| \geq 16$ and we can take $D$ in $S$. Suppose
% $D \neq S$. As $D$ has a characteristic cycle of order $|D|/2$,
% and $|x^iy| \leq 4$, that cycle is in $\langle x  \rangle$
% and thus in $\langle x^2  \rangle$. However, an element
% of $\langle x^2  \rangle$ is either fixed or 
% inverted by an element of $S$. 
 The uniqueness statement regarding the proper 
 subgroups of $T$ is an instance of 
 Corollary~\ref{conl}.  
%Let $S = \langle x, y \rangle$
% as in the statement of the lemma. If we set $X = x^2$ 
% and  $Y=xy$ then 
%$$
% \langle X, \, y \mid X^{2^{a-2}} = y^2 = 1, \, 
% X^y = X^{-1}  \rangle 
%$$
% is a normal dihedral subgroup of $S$ of order $2^{a-1}$, and
%$$
% \langle X, \, Y \mid X^{2^{a-2}} = 1, \, Y^2 = -1, \, 
% X^Y = X^{-1}  \rangle 
%$$
% is a normal generalised quaternion subgroup of $S$ of 
% order $2^{a-1}$.
%D_8: unique because any D_8 is monomial, thus conjugate
%to Sylow 2-subgroup of full monomial
\end{proof}
\begin{remark}
 It is not difficult to write down explicit generators for 
 all nonabelian 2-subgroups of $\mathrm{GL}(2,\mathbb{F})$
 from the Sylow 2-subgroup generators given in 
 Lemma~\ref{2gps}. 
\end{remark}
\begin{corollary}
\label{cclcount}
 Assuming the notation and conditions of 
 Lemma{\rm~\ref{2gps}}, the number of distinct 
 $\mathrm{GL}(2,\mathbb{F})$-conjugacy 
 classes of nonabelian nilpotent primitive  subgroups of 
 $\mathrm{GL}(2,\mathbb{F})$ is
 $r(t-1)$, where $r$ is the number of 
 positive integer divisors of $q-1$.
\end{corollary}
\begin{proof}
 If $G \leq \mathrm{GL}(2,\mathbb{F})$ is 
 nonabelian nilpotent primitive then $G=G_2 \times C$ 
 for some odd order subgroup $C$ of $\mathbb{F}^{\times}1_2$.
 %If $|G_2|$ is $8$ or $2^{a+2}$ then there is a single 
 %possibility for the isomorphism type of $G_2$, whereas 
 %there are two types otherwise, by Lemma~\ref{2gps}. 
 %existence important
 There are precisely $r/2$ choices for $C$. Using 
 Theorem~\ref{konzal} and Lemma~\ref{2gps} we then count
 $2(r/2)+2(t-2)r/2$ possible isomorphism types for $G$
 in total, so the result follows from 
 Corollary~\ref{conl}.  
\end{proof}

\section{Degree greater than 2}
\label{greater2}
 
 %We now proceed to a full characterisation
 %of nilpotent primitive linear groups over finite fields.  
 Proposition~\ref{oddfact} places a restriction on the 
 degree of a  nilpotent primitive linear group over a finite
 field. 
 We give another restriction below. The idea is to
 connect  general degree to degree 2, by means of 
 Proposition~\ref{absirredcase1} and the following.
\begin{theorem}
\label{prewehr}
 Let $G$ be primitive.
% subgroup of $\mathrm{GL}(n,\mathbb{F})$. 
%linear group of degree $n$ over $\mathbb{F}$. 
 For some $m$ dividing $n$ 
 there is an isomorphism of $G$ onto an absolutely 
 irreducible primitive subgroup of
 $\mathrm{GL}(n/m,q^m)$. 
%primitive linear group of degree 
% $n/m$ over the degree $m$ extension of $\mathbb{F}$.
\end{theorem}
\begin{proof}
 See \cite[1.19, p.12]{Wehrfritz}.
\end{proof}
%\begin{proposition}
% Suppose $G$ is nonabelian nilpotent primitive, and let 
% $2^a$ be the largest power of $2$ dividing $q+1$. 
% Then $|G_2| \leq 2^{a+2}$, and $G_2$ is semidihedral if 
% $|G_2| = 2^{a+2}$, quaternion if $|G_2| = 8$,
% and generalised quaternion or dihedral otherwise.
%\end{proposition}
%\begin{proof}
% We combine Theorems~\ref{prewehr} and \ref{konzal}
%\end{proof}
\begin{proposition}
\label{2odegree}
% If $G$ is 
 There exist  nonabelian nilpotent primitive subgroups of 
 $\mathrm{GL}(n,\mathbb{F})$ only if $n = 2m$, 
 $m$ odd. 
\end{proposition}
\begin{proof}
 By Theorem~\ref{prewehr}, suppose 
 there exist absolutely irreducible nilpotent primitive
 subgroups of $\mathrm{GL}(n/m,q^m)$, $n/m>1$.
 %$r$ dividing $n$ there is an 
 %isomorphism $f:G \rightarrow \mathrm{GL}(r, q^m)$ 
 %such that $f(G)$ is an absolutely irreducible 
 %primitive subgroup of $\mathrm{GL}(r,q^m)$, $m=n/r$.
 Proposition~\ref{absirredcase1} dictates that 
 $n/m=2$, and $q$, $q^{m}$ are both congruent to 3 mod 4 
 by Theorem~\ref{konzal}. Thus $m$ is odd.
\end{proof}

 %Perhaps the ideal situation for 
 In the sequel, our discussion of the structure of a nonabelian 
 nilpotent primitive linear group $G$
 turns on whether or not $G$ has cyclic primitive subgroups. 
 Once a cyclic subgroup $A$ of $G$ is known to be 
 irreducible, a simple order criterion 
 (Proposition~\ref{likersim}) can be used to test
 primitivity of $A$. If $A$   
 is found to be imprimitive, then it must have 
 an imprimitivity system of size 2 
 (Theorem~\ref{evenimprim}).
\begin{lemma}
\label{nakeasy}
 %Let $\mathbb{E}$ be any field, and let $W$ be the 
 %natural module for $\mathrm{GL}(n,\mathbb{E})$.
 %Suppose $H \leq \mathrm{GL}(n, \mathbb{E})$ is irreducible, 
 %with a proper subgroup $K$ of index $k$. 
%\begin{itemize}
%\item[{\rm (i)}]  The dimension of a nonzero $K$-submodule 
% of $W$ is at least $n/k$.
%\item[{\rm (ii)}]  If $W$ has a $K$-submodule $U$ of
% dimension $n/k$ then $W\cong U^H$, so that $H$
%%is the induced module $U^H$, 
 %is imprimitive. 
 Let $G$ be irreducible, with a proper subgroup $K$ 
 of index $k$. 
\begin{itemize}
\item[{\rm (i)}]  The dimension of a nonzero $K$-submodule 
  of $V$ is at least $n/k$.
 \item[{\rm (ii)}]  If $V$ has a $K$-submodule $U$ of
 dimension $n/k$ then 
%$V\cong U^G$, so that 
 $G$ is imprimitive.
\end{itemize}
\end{lemma}
\begin{proof}
 This lemma (which is true for any field in place of $\mathbb{F}$)
 follows mainly from Nakayama Reciprocity; 
% \cite[4.5, p.50]{HuppertBlackburn};
 cf. the proof of \cite[2.6]{Sim}.
% and irreducibility of $H$ yield that $V$ is isomorphic
% to a quotient of $U^H$. Thus $U^H$ has dimension at least $n$,
% and since $U^H$ has dimension precisely 
% $k \cdot \mathrm{dim} \, U$ by definition, 
% (i) follows. In case (ii), $V \cong U^H$.  
\end{proof}

 For the rest of the section, $q\equiv 3$ mod 4
 and $n=2m$, $m>1$ odd.   
\begin{theorem}
\label{evenimprim}
 Let $G= G_2 \times C$ be nonabelian nilpotent primitive, 
 and let $A$ be a cyclic subgroup of $G$ of index $2$.
 If $A$ is imprimitive then every $A$-system of  
 imprimitivity consists of $2$ components. 
\end{theorem}
\begin{proof}
 Suppose $A$ has an imprimitivity system of odd prime
 size $t$. Let $D$ be the index $t$ subgroup of $C$.
 By Proposition~\ref{likesim} the dimension of an 
 irreducible $A_2D$-submodule of $V$ is  
 $n/t$. By Clifford's Theorem $G_2D$ is completely 
 reducible, with say $l$ irreducible parts of equal degree. 
 For the same reason a $G_2D$-submodule of $V$ 
 is completely reducible as an $A_2D$-module.
 Thus $n/l$ is divisible by $n/t$, so $l$ is 1 or $t$. 
 However $l\neq t$ by 
 Lemma~\ref{nakeasy} (ii). 
 Also $l \neq 1$, for if $G_2D$ were irreducible then 
 an $A_2D$-submodule of $V$ would have dimension at
 least $m$ by Lemma~\ref{nakeasy} (i).
 Thus the only prime dividing the size $s$ of an
 $A$-system of imprimitivity is 2 by Lemma~\ref{sizeprim},
 and as $s$ divides $2m$, we are done.
\end{proof}

 Until further notice, we fix the following notation and 
 hypotheses. Let $G = G_2 \times C$ where 
 $C = \langle c\rangle$ is odd order cyclic, 
 and $G_2$ is (nonabelian) generalised quaternion, 
 dihedral or semidihedral. 
 That is, $G_2 = \langle A_2, g\rangle$, 
 $|G_2 : A_2| = 2$, 
 %$g^2$ $\in$ $A_2$ is central of order no more than 2, 
 %$\langle -1_n \rangle \leq A_2$, and 
 and $A_2 = \langle a\rangle$, $|a| \geq 4$. 
 Suppose $A = \langle c, a\rangle$ is irreducible. 
 Then $g^2$ $\in$ $\langle -1_n \rangle \leq A_2$.
 We define a field automorphism $\delta$ of 
 $\Delta  = \langle A\rangle_\mathbb{F}$ by
 $\delta: x \mapsto g^{-1}xg$, 
 $x \in \Delta$. Let $\mathbb{K}$ be the $\delta$-invariant 
 subfield  $\{x \in \Delta: gx = xg\}$ of $\Delta$. 
 Note that $C \subseteq \mathbb{K}$. In addition
 $|\Delta : \mathbb{K}|$ $=$ $2$, so
%Isaacs Algebra, p.285, 18.20
 $|\mathbb{K} : \mathbb{F}1_n| = m$ and thus 
 $|\mathbb{K}|$ $\equiv$ 3 mod 4.
%\begin{lemma}
% $\mathbb{K}$ is the centraliser of $G$ in 
% $\mathrm{Mat}(n, \mathbb{F})$.
%\end{lemma}
%\begin{proof}
% Let $P$ be the centraliser of $G= \langle A, g\rangle$ 
% in $\mathrm{Mat}(n, \mathbb{F})$.
% By definition $\mathbb{K} \subseteq P$.
% Since $P$ centralises $A$, by Schur's Lemma
% every nonzero element of $P$ is in 
% $\mathrm{GL}(n,\mathbb{F})$. Thus $P\subseteq \Delta$, 
% because $\Delta^{\! \times}$ is the
% centraliser of $A$ in $\mathrm{GL}(n,\mathbb{F})$, so
% $P \subseteq \mathbb{K}$. 
%\end{proof}
\begin{remark}
 %It can be shown that 
 $\mathbb{K}$ is the centraliser of $G$ in 
 $\mathrm{Mat}(n, \mathbb{F})$.
\end{remark}
\begin{lemma}
\label{delinv}
 $|\langle A_2\rangle_\mathbb{F} : \mathbb{F}1_n|=2$ 
 and $\langle A_2\rangle_\mathbb{F} \cap \mathbb{K}$ 
 $=$ $\mathbb{F}1_n$.
\end{lemma}
\begin{proof}
 Let $B$ be the subgroup of $A_2$ of order 4. Since 
 $B^2$ but not $B$ is scalar,  
 $|\langle B \rangle_\mathbb{F} : \mathbb{F}1_n| =2$.  
 %If $|a_2| >4$ then let $B'$ be the subgroup of $A_2$ 
 %containing $B$ maximally. Now  
 %$|\langle B' \rangle_\mathbb{F}$ has an
 %$\mathbb{F}$-spanning set of size 4 and contains a 
 %subfield of $\mathbb{F}$-dimension 2, so 
 %$|\langle B' \rangle_\mathbb{F}  : \mathbb{F}1_n| =2$
 %(can't be a power of 2 any bigger than 2).  
 An inductive argument then shows that 
 $|\langle A_2\rangle_\mathbb{F} : \mathbb{F}1_n| =2$.
 For the second claim we need only  
 observe that $\langle A_2 \rangle_\mathbb{F}$   
 contains an element of order 4 while $\mathbb{K}$ 
 does not.
\end{proof}

We denote a diagonal matrix as the ordered tuple of its main
diagonal entries.
\begin{lemma}
\label{prea2c}
% Let $\mathbb{E}$ be any field, $B$ an abelian
% irreducible  subgroup of $\mathrm{GL}(r,\mathbb{E})$,
% and $\bar{B}\leq B$. The irreducible $\bar{B}$-submodules 
% of the natural module $W$ for $\mathrm{GL}(r,\mathbb{E})$
% have dimension $|\langle \bar{B} \hspace*{.2mm} 
% \rangle_{\mathbb{E}} : \mathbb{E}1_r|$, 
% and they are all isomorphic to each other.
If $\bar{A}\leq A$ then the irreducible $\bar{A}$-submodules
 of $V$ have dimension $|\langle \bar{A} \hspace*{.2mm} 
 \rangle_{\mathbb{F}} : \mathbb{F}1_n|$, 
and they are all isomorphic to each other.
\end{lemma}
\begin{proof}
% \cite[Theorem 4, p.99]{Suprunenko}.
 By Clifford's Theorem,  
 $\bar{A}$ is conjugate to  
 $\langle (x, \ldots \, \! , x) \rangle$ 
 for some irreducible subgroup $\langle x \rangle$
 of $\mathrm{GL}(r, \mathbb{F})$,
 $r=|\langle x\rangle_{\mathbb{F}}:\mathbb{F}1_{r}|$ 
 dividing $n$. 
 Clearly then $\langle \bar{A} \hspace*{.2mm} 
 \rangle_{\mathbb{F}}$ $\cong$ 
 $\langle  x\rangle_{\mathbb{F}}$.    
\end{proof}
\begin{lemma} 
\label{a2cirred}
\begin{itemize}
\item[{\rm (i)}] 
 An irreducible $A_2$-submodule of $V$ has
 dimension $2$.
\item[{\rm (ii)}] 
 $\mathbb{K} = \langle C\rangle_\mathbb{F}$, and an 
 irreducible $C$-submodule of $V$ has dimension 
 $m$.
\end{itemize}
\end{lemma}
\begin{proof}
%The compositum of $\mathbb{K}$ and 
% $\langle A_2\rangle_{\mathbb{F}}$ in $\Delta$ contains $A$ and  
% is therefore $\Delta$. Clearly 
 Lemmas~\ref{delinv} and \ref{prea2c} give (i).
 Since $\Delta$ is the  compositum of 
 $\langle C\rangle_{\mathbb{F}}$ and  
 $\langle A_2\rangle_{\mathbb{F}}$,  
 $\langle C\rangle_\mathbb{F}\subseteq \mathbb{K}$, and
 $\langle A_2\rangle_{\mathbb{F}} \cap \mathbb{K}$ 
 $=$ $\mathbb{F}1_n$, we get that
 $|\langle C \rangle_\mathbb{F}  : \mathbb{F}1_n| =m$
 i.e. $\langle C\rangle_\mathbb{F}=\mathbb{K}$,
 as required for (ii). 
%This degree is $m$ if $\bar{A}=C$, and $2$ if $\bar{A}=A_2$, so 
\end{proof}
%\begin{remark}
% By Corollary~\ref{a2cirred}, $a_2$ is conjugate to an
% $n/2\times n/2$ block scalar matrix, where the constant
% $2 \times 2$ block generates an irreducible abelian
% subgroup of $\mathrm{GL}(2, \mathbb{F})$. Therefore
% $|a_2|$ divides $2(q+1)$. A similar analysis is possible on
% $C$; its order divides $q^{n/2}-1$ but not $q^t-1$ for
% any $t$ properly dividing $n/2$.   
%\end{remark}
%
%\begin{lemma}
%\label{no7} 
% Let $C = \langle (c_1, c_1) \rangle$ where
% $\langle c_1\rangle \leq  \mathrm{GL}(m, \mathbb{F})$
% is primitive. Suppose $V$ is the direct sum of  
% $C$-submodules $W_1$ and $W_2$. Then $W_1, W_2$ are 
% isomorphic irreducible $C$-modules 
% and $C\! \mid_{W_i}$ is primitive.
%\end{lemma}
%\begin{proof}
% By the Jordan-H\"{o}lder Theorem, $W_1$ $\cong$
% $W_2$ and $\bar{C} =C\! \! \mid_{W_1}$ is an
% irreducible subgroup of $\mathrm{GL}(W_1)$. Since 
% $|\bar{C}| = |c_1|$,  
% $\bar{C}$ is primitive by Proposition~\ref{likersim}. 
%\end{proof}
%
\begin{theorem}
\label{firstclassif}
% Let $G = G_2 \times C \leq  \mathrm{GL}(2m, \mathbb{F})$, 
% where $C =\langle c\rangle$ is odd order cyclic, 
% $G_2 = \langle a, g\rangle$ is the Sylow 
% $2$-subgroup of $G$, $A_2 = \langle a\rangle$ is a maximal 
% abelian subgroup of $G_2$ of index $2$, 
% $A = A_2 \times C$
%% = \langle c, a_2\rangle$ 
% is irreducible, and  
% $|\langle C\rangle_\mathbb{F} : \mathbb{F}1_n| = m$. 
% Assume the notation and hypotheses introduced before 
% Lemma$~\ref{}$.
 $G$ but not $A$ is primitive if and only if $A$ has an 
 imprimitivity system 
%$\mathcal{S} = \{ V_1 ,  V_2\}$,
 of size $2$,  
 $C$ is conjugate to $\langle (x,x) \rangle$ 
%$(C_1, C_1)$ 
 for some primitive subgroup 
%$C_1= 
$\langle x \rangle$ of 
 $\mathrm{GL}(m, \mathbb{F})$, and 
 $a^2 = g^2 = [a, g] = -1_n$.
\end{theorem}
\begin{proof}
% We note at the outset that if $K$ is any subgroup of
% $A$ then $\langle K \rangle_{\mathbb{F}}$ is a field.
 Suppose $G$ but not $A$ is primitive.
 By Theorem~\ref{evenimprim},
 $A$ has an imprimitivity system $\mathcal{S} = \{ V_1, V_2 \}$.
% By irreducibility of $A$, $\mathrm{dim}\,  V_1$ $=$
% $\mathrm{dim}\,  V_2=m$.
% Recall the notation 
 Let $\phi$ be the permutation representation of $A$ in 
 $\mathrm{Sym}(2)$ arising from the action on $\mathcal{S}$, 
 and set $\ker \phi=A'$.
 Since $|C|$ is odd, necessarily $C \leq A'$, 
and then by
%  $C$ $=$
 %$\mathrm{diag} (C_1, C_2)$, $C_i \leq \mathrm{GL}(m, \mathbb{F})$, 
% $\{ \mathrm{diag} (\rho_1(c), \rho_2(c)) \mid c \in C \}$, 
% $C_i \leq \mathrm{GL}(m, \mathbb{F})$, 
% up to conjugacy. By 
 Lemmas~\ref{prea2c} and \ref{a2cirred},
%, $V_1$ and
% $V_2$ are irreducible $C$-modules, and because
% (clearly) $V_1 \cong V_2$ as $C$-modules, in fact 
% $C_1$ is $\mathrm{GL}(m, \mathbb{F})$-conjugate to $C_2$. 
% Hence 
% $C_1 = C_2$ is irreducible, and we may take 
% $x=y$. 
 $C = \langle (x,x) \rangle$ up to conjugacy,
 where $C_1 = \langle x \rangle \leq \mathrm{GL}(m, \mathbb{F})$ 
 is irreducible. Since $\langle C\rangle_\mathbb{F}$ $\leq$
 $\langle A'\rangle_\mathbb{F}$
 and $|\Delta : \langle C\rangle_\mathbb{F}| =2$, 
%$$
%2m = |\Delta: \mathbb{F}1_n| = 
%  |\Delta : \langle C\rangle_\mathbb{F}|
%  |\langle C\rangle_\mathbb{F} : \mathbb{F}1_n|  
%$$ 
 either $\langle A' \rangle_\mathbb{F}$ $=$ 
 $\Delta$ or $\langle A'\rangle_\mathbb{F} = 
 \langle C\rangle_\mathbb{F}$. 
 If $\langle A' \rangle_\mathbb{F} = 
 \Delta$ then $A'$ is irreducible, 
 which is definitely not true.
 Thus $\langle A'\rangle_\mathbb{F} = 
 \langle C\rangle_\mathbb{F}$.    
 An element $h\neq 1$ of $A' \cap A_2$ has 2-power order;
 but $h$ $\in$  $\langle C\rangle_\mathbb{F}$, which
 does not contain an element of order 4, so $h = -1_n$. 
 %$|h| = 2$. 
 %%We have that $h \in \langle (x, y)\rangle$
 %Write $h = (x,x)$ for some 
 %$x \in \langle C_1 \rangle_{\mathbb{F}}$, $|x| = 2$. 
 %%$\langle C_1 \rangle_{\mathbb{F}}$ contains all scalars
 %%of $\mathrm{GL}(m,\mathbb{F})$, so the element of 
 %%order 2 in  $\langle C_1 \rangle_{\mathbb{F}}$ is 
 %In fact $x =  -1_m$ because $C_1$ 
 %is irreducible.
 %singer cycle contains all scalars
 %Consequently $h \in \langle -1_n\rangle$, i.e.
 %$A_2 \cap A' \leq \langle -1_n\rangle$. 
 Consequently 
 $A' = (A' \cap A_2)\times C = \langle -1_n, C \rangle$ and 
 $A'|_{V_i} = \langle -1_m, C_1 \rangle$. Since 
 $A'|_{V_i}$ is primitive (otherwise we get an $A$-system of
 imprimitivity of size greater than 2, violating 
 Theorem~\ref{evenimprim}), $C_1$ is primitive.
 From $|A: A'|=2$ and $|A'| = 2|C|$ we infer $|A_2| =4$
 and $|G_2|=8$. 
 By Theorem~\ref{konzal}, then, $G_2 \cong Q_8$.
 %($D_8$ has a noncyclic abelian normal subgroup, 
 %which is normal in $G$). Note that $-1_n \in G_2$.
 We have seen that $|a|=4$, and 
 $g \notin Z(G_2) = \langle -1_n \rangle$
 implies 
%$|g| = 4$. Hence 
 $a^2=g^2= [a,g] = -1_n$.

 %Now suppose $A$ has imprimitivity system $\mathcal{S}$,
 %$C$ and $G_2 \cong Q_8$ are as described and 
 Now we prove the converse. Suppose 
 %$A$ has imprimitivity system $\mathcal{S}$, 
 $C = \langle (x,x) \rangle$ where 
 % (C_1, C_1)$ for a 
 $\langle x \rangle := C_1$ is a primitive subgroup of 
 $\mathrm{GL}(m, \mathbb{F})$, and $G_2 \cong Q_8$.
% and $-1_n \in G_2$. 
 Also suppose $\mathcal{T}$ $=$ $\{  W_1 , \ldots \, \! ,  W_k\}$
 is a $G$-system of imprimitivity, and let $\varphi
 : G \rightarrow \mathrm{Sym}(k)$ be the associated 
%transitive 
 permutation representation.
 Of course, $\mathcal{T}$ is simultaneously an $A$-system of 
 imprimitivity, so that
% for every subgroup of $G$, and in particular 
 if $k$ is divisible by an odd prime $t$
 then by Lemma~\ref{sizeprim}, $A$ has an imprimitivity 
 system of size $t$.
 By Proposition~\ref{likesim}, $|A| = 4|C_1|$ divides 
 $t(q^{m/t}-1)(q^{m/t}+1)$.  
 Some prime divisor $r$ of $|C_1|$ must divide $q^{m/t}+1$
 by Proposition~\ref{likersim}, and therefore $r$ divides $q^m+1$.
 %As $C_1$ $\leq$ $\mathrm{GL}(m,\mathbb{F})$ is irreducible,
 However $|C_1|$ divides $q^m-1$, so $r$ divides $2q^m$, 
 a contradiction. 
%It follows that $k$ is a power of 2. 
% Since $k$ divides $2m$, $k=2$.
 Hence $k=2$.
% By Lemma~\ref{no7}, 
% 
%As before it can be shown that 
 %$\langle A \cap \ker \varphi \rangle_{\mathbb{F}}
 %= \langle C \rangle_{\mathbb{F}}$, which has 
 %no element of  
%
%is the $\mathbb{F}$-enveloping 
% algebra of $A$ $\cap$ $\ker \varphi$. 
%
% order 4, and thus $A_2 \cap \ker \varphi$ $=$
% $\langle -1_n\rangle$. 
%
% (which is a subgroup of $G_2$ by hypothesis). 
 
 Since
% $\varphi(C)$ is trivial, 
 $\varphi(A_2)$ is transitive, after replacing $G$ 
 with a suitable conjugate we have $a = (u, v) w$
 where $u, v$ $\in$ $\mathrm{GL}(m, \mathbb{F})$ and
$$
 w = 
 \left( \begin{array}{cc}
 0_m & 1_m \\
  1_m & 0_m 
 \end{array}
 \right) ,
$$
 %and $C\leq A'$ is $2\times 2$ block diagonal
 %form, with the $m\times m$ blocks conjugate 
 %in $\mathrm{GL}(m,\mathbb{F})$. A further conjugation leaves
 %$a$ in the same form, 
 with $C$ retaining its original form
 %$\mathrm{diag}(C_1,C_1)$. 
 $\langle (x,x) \rangle$. The relation 
 $a^2 = -1_n$ yields $v$ $=$ $-u ^{-1}$.
% $a = (u, -u^{-1})T_2$. 

 Recall Proposition~\ref{normcentr}.
 Since $a$ centralises $C$, $u$ centralises $C_1$,
 and therefore $u \in \langle C_1\rangle_\mathbb{F}$.
% $\mathrm{GL}(m,\mathbb{F})$-centraliser of the 
% irreducible subgroup $C_1$ of  $\mathrm{GL}(m,\mathbb{F})$,
% viz. $\langle C_1\rangle_\mathbb{F}$. Since
% As $g$ normalises $A$, it normalises 
% $\Delta^{\! \times}$.
%= \langle C, a\rangle_\mathbb{F}$.
 Note that $d =  (1_m, -1_m)$ inverts $a$, so
  it normalises but does not centralise 
 $\Delta^{\! \times}$. Certainly $g$ normalises  
 $\Delta^{\! \times}$, so
%$e \not \in \Delta$, 
  $g = db$ for some $b \in \Delta^{\! \times}$.
% by Proposition~\ref{normcentr}. 
%Since $\mathcal{T}$ is a $G$-system of imprimitivity, either
 
 Over $\langle C \rangle_\mathbb{F}$, $\Delta$ has basis
 $\{ 1_n, a\}$.
 Thus if 
%$h \in \Delta$ acts trivially on $\mathcal{T}$ then 
%$h \in \langle C \rangle_\mathbb{F}$.
% 
% Either $g \in \ker \varphi$ or $|\varphi (g)| = 2$.
 $g \in \ker \varphi$ then $b$ $=$ 
 $(y, y)$ for some
 %$b$ $=$ 
 %$(v, v) \in \langle C\rangle_\mathbb{F}$, 
 %i.e. 
 $y\in \langle C_1\rangle_\mathbb{F}$. Since $g^2 = -1_n$
 it follows that $y^2 = -1_m$,
% That is, $|y| =4$, 
 which is impossible.

 Suppose $g \notin \ker \varphi$, so that
% i.e. $\langle g\rangle$ acts on $\mathcal{T}$ transitively. 
 %$e$ fixes each of $W_1$ and $W_2$, 
 %$b \in \Delta$ must interchange them, i.e. 
 $b$ $=$ $(u y ,-u^{-1} y)w$ for some 
 $y \in \langle C_1\rangle_\mathbb{F}$.
% so $g$ $=$ $(u x ,u^{-1} x)w$.
 %Then $g$ $=$ $eb$ $=$ 
 %$(v u , v u^{-1})T_2$. Since 
 Once again, $g^2 = -1_n$ implies 
 $\langle C_1\rangle_\mathbb{F}$ has an element
% (here, $u$)???
 of order 4. This final contradiction proves that $G$ is primitive. 
\end{proof}

 We next give an explicit description of 
 the groups in Theorem~\ref{firstclassif},
 which can
 be extracted from the proof of the theorem.
% and also
% serves as an existence statement about the groups.  
 An auxiliary fact is needed: for any prime 
 $p$, $-1$ is a sum of two squares mod $p$.
%clear for p=2 (e=1, f=0) and p equiv 1 mod 4 (e=i, f=0)
%Suppose p equiv 3 mod 4. Then -1 is a representative
%of Z_p*/Z_p*^2, so either we find an e,f, or Z_p*^2
%is closed under addition. But then Z_p*^2 U 0 is a subfield
%of Z_p, so must have char p and be Z_p, contradiction
\begin{proposition}
 Let $C_1=\langle x \rangle$ be a cyclic primitive subgroup of 
 $\mathrm{GL}(m, \mathbb{F})$, and choose 
 $e , f \in  \langle C_1\rangle_\mathbb{F}$ such that
 $e^2+f^2=-1_m$. 
 Define $C = \langle (x, x)\rangle$ and 
 $G_2 = \langle a, g\rangle$, 
 where
$$a =\left( \begin{array}{rr}
0_m & 1_m \\
-1_m & 0_m 
\end{array} \right) , \ \ \ g = \left( \begin{array}{lrc}
e & f \\
f & -e 
\end{array} \right) .
$$
 Then $G = G_2 \times C$ is a nonabelian nilpotent 
 primitive subgroup of $\mathrm{GL}(n, \mathbb{F})$
 with an abelian imprimitive subgroup of index $2$.
 Moreover, any group of this kind with the 
 same order as $G$ is conjugate to $G$.
\end{proposition}
 %To verify these assertions, first note from the proof
 %of Theorem~\ref{firstclassif} that $C$ is as stated above, 
 %and $a = (u, -u^{-1})T_2$,
 %$u \in \langle C_1 \rangle_{\mathbb{F}}$. Then 
 %$a$ is conjugate by $x= (u, 1_m)$ to
 %$(1_m, -1_m)T_2$. Now $C^x = C$ and $g^x$ 
 %satisfies $(g^x)^2=  [a^x, g^x] = -1_n$, so as in the 
 %proof of Theorem~\ref{firstclassif} after relabelling 
 %$g^x$ as $g$ we get that 
 %$g=eb$ where $e =(1_m, -1_m)$ and
 %$b \in \Delta$.
 %We know $\Delta$ has basis $\{ 1_n, a\}$ over
 %$\langle C\rangle_{\mathbb{F}}$, and so can work
%out the form of an element of 
 %$\Delta$, and then see that $g$ has the form
%$$
%\left( 
%\begin{array}{cc}
% \gamma & u\beta \\
% -u^-1\beta & -\gamma
%\end{array}
%\right) ,  \ \ \  \beta, \gamma 
% \in \langle C_1 \rangle_{\mathbb{F}}
%$$
%Conjugate by $(u, 1_m)$
% Enforcing the relation $g^2 = -1_n$ yields  
% $\gamma^2+\beta^2 = -1_m$, and we are done
%% $\gamma \beta \alpha$ $=$  $\gamma \beta \alpha^{-1}$,
%% $\gamma^2+\alpha^2\beta^2 = 1_m$, and  
%% $-\gamma^2-\alpha^{-2}\beta^2 = -1_m$.
%% If $\gamma=0_m$ then $\beta^2 = -1_m$, but 
%% $\langle C_1 \rangle_{\mathbb{F}}$ has no element 
%% of order 4. Similarly we cannot have $\beta = 0_m$.
%% Thus $\alpha = \alpha^{-1}$, i.e. $\alpha = -1_m$.
%% Relabelling $\gamma$ as $e$ and $-\beta$ as $f$ then 
%% validates the above assertions.

 Our last undertaking in this section is to extend 
 Corollary~\ref{conl}.
% to higher degrees.
% can say a little more about the 
% $\mathrm{GL}(n,\mathbb{F})$-conjugacy classes of 
% the groups in case 3 above.
%, and for this we need the following facts about the behaviour
% of irreducible representations under extensions
% of the ground field.
 We use the following version 
 of Theorem~\ref{prewehr}.
\begin{theorem}
\label{explwehr}
 Let $H$ be an irreducible but not absolutely 
 irreducible subgroup of ${\rm GL}(r,\mathbb{F})$, $r>1$. 
 For some $s>1$ dividing $r$ there is an 
 absolutely irreducible subgroup $\hat{H}\cong H$ of
 ${\rm GL}(r/s,q^s)$ such that $H$ is
 ${\rm GL}(r, q^s)$-conjugate to the 
 group $\tilde{H}$ of block diagonal matrices
$$
(h, h^{\sigma}, \ldots \, \! , 
h^{\sigma^{s-1}}) , \ \ \ \ h \in \hat{H} , 
$$
 where $\sigma$ is 
%the Galois automorphism of 
 a fixed generator 
% $\mathrm{GF}(q^m)$ defined by $f \mapsto f^q$, 
 of $\mathrm{Gal}(\mathrm{GF}(q^s)/\mathbb{F})$, 
 and $h^{\sigma^j}$ denotes the matrix resulting from 
 entrywise action of $\sigma^j$ on $h$. 
\end{theorem}
\begin{proof}
 See \cite[9.21, p.154]{Isaacs} and the proof of that 
 result.
\end{proof}
%\begin{remark}
 %%As observed in the proof of Proposition~\ref{2odegree},
 %%$G_2$ is isomorphic to a 2-subgroup of 
 %%$\mathrm{GL}(2,q^m)$, 
 %The Sylow 2-subgroups of $\mathrm{GL}(2,\mathbb{F})$
 %and of $\mathrm{GL}(2,q^m)$ coincide. 
 %Theorem~\ref{explwehr} and Lemma~\ref{2odegree}
 %allow us to display 
 %%so $|G_2|$ divides $4(q+1)$. 
%\end{remark}
\begin{theorem}
\label{iecagain}
 Nilpotent primitive subgroups $G$ and $H$ of 
 $\mathrm{GL}(n,\mathbb{F})$ 
 are conjugate in  $\mathrm{GL}(n,\mathbb{F})$
 if and only if they are isomorphic.
 \end{theorem}
\begin{proof}
 Suppose $G \cong H$.
% and $n>2$. 
 Since $G$ and $H$ are not
 absolutely irreducible ($n>2$) by Proposition~\ref{absirredcase1},
 there are absolutely irreducible 
 subgroups $\hat{G}$, $\hat{H}$ of $\mathrm{GL}(2,q^m)$, and 
 subgroups $\tilde{G}$, $\tilde{H}$ of $\mathrm{GL}(n,q^m)$
 constructed from  $\hat{G}$, $\hat{H}$ as in 
 Theorem~\ref{explwehr}, such that $G$ is conjugate to 
 $\tilde{G}$ and $H$ is conjugate to 
 $\tilde{H}$. The groups $\hat{G}$ and $\hat{H}$ are primitive, 
 %really? can dispense with Q_8 by previous result, then orders
 %show hatG,hatH primitive in GL(2)
 so they are $\mathrm{GL}(2,q^m)$-conjugate by Corollary~\ref{conl}.
 Then it is easy to see that $\tilde{G}$ and $\tilde{H}$
 are $\mathrm{GL}(n,q^m)$-conjugate. 
 %conj by (x, x^sigma, x^sigma^2, ...
 Hence $G$ and $H$ are $\mathrm{GL}(n,\mathbb{F})$-conjugate 
 by the Deuring-Noether Theorem.
\end{proof}

\section{Summary and concluding remarks}
\label{summary}

 Suppose $q\equiv 3$ mod 4. Let $G$ be nonabelian 
 nilpotent primitive, and let $C$ be the 
 direct product of the odd order Sylow subgroups of 
 $G$. Previously we have established that one of 
 the following situations must occur.
\begin{enumerate}
\item $n=2$, $G_2$ is absolutely irreducible, 
 and $C$ is scalar.
\item $n=2m$, $m>1$ odd, $G_2 \cong Q_8$, and
 $G$ has an imprimitive cyclic subgroup of index 2. 
\item $n=2m$, $m>1$ odd, and
 $G$ has a primitive cyclic subgroup of index 2. 
\end{enumerate}
 In all cases, $G_2$ is either $Q_8$,
  or a Sylow 2-subgroup of $\mathrm{GL}(n,\mathbb{F})$,
  or is generalised  quaternion or dihedral of order at least
 16. 
%
% In case 1 above, two nilpotent primitive subgroups of
% $\mathrm{GL}(2,\mathbb{F})$ are conjugate if and 
% only if they are abstractly isomorphic, by Lemma~\ref{2gps}.
% In case 1, all subgroups of $G$ of index 2 are primitive???
%NO
 

 Suppose $G \cong Q_8 \times C$
 and $n>2$. 
%By Theorem~\ref{prewehr} and Proposition~\ref{absirredcase1}, 
 Then $|C|$ divides $q^m-1$,
 so an index 2 subgroup of $G$ has order dividing 
 $2(q^m-1)$. Such a subgroup of 
 $G$ is imprimitive by Proposition~\ref{likersim}.
 Hence $|G_2| > 8$ in case 3, 
% A group of the second type above cannot be conjugate
% to a group of the third type, as may be verified
% %For if $H\neq G$ is an irreducible subgroup of $G$ in case 2 
% %then $|H \cap G_2|>2$, so $H$ is cyclic of index 2 in $G$ 
% %and thus $H$ is imprimitive 
% by Proposition~\ref{likersim}.
 %($H$ has the same order as an imprimitive subgroup of 
 %$G$). 
 %
% As we proved in 
% Section~\ref{greater2},
 and for each possible value of $|C|$, i.e. 
 the order of an odd order cyclic primitive subgroup 
 of $\mathrm{GL}(m, \mathbb{F})$, 
%positive integer $r$ dividing $q^m-1$ but 
% not $k(q^{m/k}-1)$ for any prime divisor $k$ of $m$,
 there is a single conjugacy class of primitive 
 subgroups of $\mathrm{GL}(n,\mathbb{F})$ isomorphic
 to $Q_8 \times C$. The $\mathrm{GL}(n,\mathbb{F})$-conjugacy 
 classes of the groups in case 3 may be counted with the
 aid of Theorem~\ref{iecagain} (cf. Corollary~\ref{cclcount}).  

\subsection*{Acknowledgment}
% \mbox{A.} \mbox{S.} Detinko
 The first author received support 
 from the Enterprise Ireland International Collaboration 
 Programme, and is also indebted
 to the Mathematics Department at NUI, Galway, where
 work on this paper was carried out.

\bibliographystyle{amsplain}
\begin{thebibliography}{10}
\bibitem{Conlon}
S.~B. Conlon, \emph{$p$-groups with an abelian maximal subgroup and cyclic
  center}, J. Austral. Math. Soc. \textbf{22} (1976), 221--233.

\bibitem{Dixon}
John~D. Dixon, \emph{The structure of linear groups}, Van Nostrand Reinhold,
  London, 1971.

\bibitem{HuppertBlackburn}
B.~Huppert and N.~Blackburn, \emph{Finite groups {II}}, Springer-Verlag,
  Berlin, 1982.

\bibitem{Isaacs}
I.~Martin Isaacs, \emph{Character theory of finite groups}, Dover, New York,
  1994.

\bibitem{Konyuh}
V.~S. Konyukh, \emph{Locally nilpotent linear groups}, Dokl. Akad. Nauk BSSR
  \textbf{28} (1984), no.~3, 197--199 (Russian).

\bibitem{LGP}
C.~R. Leedham-Green and W.~Plesken, \emph{Some remarks on {Sylow} subgroups of
  general linear groups}, Math. Z. \textbf{191} (1986), 529--535.

\bibitem{Sim}
Hyo-Seob Sim, \emph{Metacyclic primitive linear groups}, Comm. Algebra
  \textbf{22} (1994), no.~1, 269--278.

\bibitem{Suprunenko}
D.~A. Suprunenko, \emph{Matrix groups}, Transl. Math. Monogr., vol. 45,
  American Mathematical Society, Providence, RI, 1976.

\bibitem{Wehrfritz}
B.~A.~F. Wehrfritz, \emph{Infinite linear groups}, Springer-Verlag, Berlin,
  Heidelberg, New York, 1973.

\bibitem{Zalesskii}
A.~E. Zalesskii, \emph{Supersolvable and nilpotent subgroups of simple
  algebras}, Dokl. Akad. Nauk. BSSR \textbf{7} (1963), no.~12, 800--802
  (Russian).
\end{thebibliography}



\end{document}
